Web agents often fail when a website changes in ways that preserve the user's goal but disrupt learned interface shortcuts. We present \webcoevo{}, a single-site but end-to-end auditable loop for generating capability-targeted website drift and testing compact reflection rules. The loop profiles failures, builds fresh \linkding{} website overlays, mines matched trajectory transitions, edits a small reflection rulebook, and evaluates the resulting policy with prompt-injection evidence recorded in trace JSONL. The upgraded WebCoEvo evidence covers three website versions: \linkding{} V1.45.0, a moderately modified website V2, and a harder generated website suite V3 over seven drift categories. On the 68 training tasks from \linkding{} V1.45.0, \expel{}-style rules raise Qwen3-VL success from \(9/68=13.2\%\) to \(56/68=82.4\%\). Under website V2, R1 reflection rules improve over \expel{}-style prompting from \(60/68=88.2\%\) to \(67/68=98.5\%\) on the training tasks and from \(114/167=68.3\%\) to \(143/167=85.6\%\) on the held-out validation tasks. On website V3, \expel{}-style prompting collapses to \(8/68=11.8\%\) on the training tasks and \(66/167=39.5\%\) on the held-out validation tasks, while R1 reaches \(65/68=95.6\%\) and \(97/167=58.1\%\). A GPT-5.4-assisted R2 update ties R1 on the held-out validation tasks but does not improve the training tasks, showing that iterative reflection is useful but not monotonic. We claim auditable \linkding{} robustness improvement from compact reflection rules, not universal web-agent robustness.
